% ============================================================
% 06-discussion.tex — V15 Discussion
% ============================================================
\section{Discussion}
\label{sec:discussion}

%------------------------------------------------
\subsection{Architectural Findings and Operational Trade-offs}

\textbf{RQ1 (Architecture-Level Effectiveness):}
The tested orchestration bundle achieves a higher Automated Fact Pass Rate point estimate with fewer specialist input tokens than Single Agent. Because the Fact Coverage confidence interval spans zero and latency measures different pipeline stages, neither completeness nor speed advantages follow. The 50\% threshold binarizes coverage and does not establish more complete answers. The No Specialist Policy ablation is consistent with a contribution from domain instructions under evaluated conditions.

\textbf{RQ2 \& RQ3 (Mechanisms, Tool Isolation, and Registry Scaling):}
The ablations associate specialist policies and heterogeneous evidence access with larger pass-rate losses when removed. Disabling deterministic calculation yields an estimated $-4.0$\,pp point change in overall automated fact pass (95\% CI $[-9.0, 0.0]$\,pp); without route repair, cross-domain pass falls from 80.0\% to 65.0\%. At 11 tools, removing isolation increases input tokens without changing the pass estimate. At 51 tools, Routed Specialist records 88.0\% tool E2E Success, $9.3$\,pp above Global Top-10 (CI $[+2.0, +17.3]$\,pp); its $4.7$\,pp difference over Full-Registry Single Agent remains uncertain (CI $[-1.3, +11.3]$\,pp). These results suggest a benefit from bounded tool exposure under tested ambiguity.

\textbf{Operational Trade-offs and Deployment Guidance:}
Single-agent timing starts post-retrieval, whereas multi-agent figures include routing; these values cannot rank architectures by speed. Specialist execution records $1{,}731$\,ms and $677$ output tokens more than generic execution. A single agent suits narrow tasks with small tool pools; scoped specialists suit institutions with distinct evidence formats and fee calculations. Institutional transfer requires mapping curriculum schemas, populating graph/vector stores, encoding fee rules in tools, and calibrating intent prompts against local regulations.

%------------------------------------------------
\subsection{Relation to Prior Work}

ReAct~\cite{yao2023react} and Toolformer~\cite{schick2023toolformer} study tool use; AutoGen~\cite{wu2023autogen} supports agent coordination. \system{} instead routes to bounded specialists with evidence paths chosen by data modality, differing from Adaptive-RAG's complexity-based retrieval~\cite{jeong2024adaptiverag}. The Source Recall drop under Uniform Text-RAG aligns with GraphRAG's rationale for relational evidence~\cite{edge2024graphrag}. These frameworks were not directly compared in our experiments.

%------------------------------------------------
\subsection{Threats to Validity}
\label{sec:limitations}
Pipeline latency ($10.65$\,s) limits interactive use. Because Single Agent timing begins post-retrieval, measurements cannot rank cross-architecture speed. Queries cover one institution; Scenario~3 uses synthetic tool neighbors. Transfer to other institutions and fee rules remains untested. Excluding Redis history leaves multi-turn dynamics unmeasured. Substring matching can pass an incorrect financial answer: \texttt{HOUT-DIR-FIN-03} reached 50\% coverage despite a 192.183M versus 40M VND discrepancy. Tool E2E checks verify execution validity rather than numeric exact match. Query family subsets have modest sample sizes ($n=19$--$33$) despite 10,000 bootstrap resamples.
